% profiles/louisiana-state-university-and-agricultural-and-mechanical-college.tex
% Louisiana State University and Agricultural & Mechanical College
% ============================================================================
% source: https://www.lsu.edu/graduateschool/current_students/theses_and_dissertations/index.php
%         (Graduate School, "Theses & Dissertations")
% source: https://www.lsu.edu/graduateschool/students/files/thesisanddissertationhandbook_feb2026.pdf
%         ("Thesis & Dissertation Handbook", February 2026 edition; the sample
%         title page is its page 25)
% accessed: 2026-09-13
% checked:  texlib-thesis-profile-round
%
% TWO EDITIONS ARE LINKED FROM THE SAME PAGE. Do not swap the source above for
% thesisanddissertationhandbook_2025.pdf, which is also live on lsu.edu and is
% reachable from a sidebar card on that page. The February 2026 edition is the
% one the page's own prose points at — "Early in your writing process, review
% the Thesis & Dissertation Handbook. It contains the complete formatting
% guidelines" — and the one its Quick Links label as "With Formatting
% Guidelines". A third file, thesisanddissertationhandbook_fall2022.pdf, still
% answers on a older path and turns up in search results.
%
% VERIFIED and set here: geometry, spacing, title page, the absence of an
% approval page, and LSU's published accessibility rule.
% NOT set: chapter-opening margin (LSU states none, and forbids exceptions to
% the margin generally) and the typeface (LSU names no face).

\thesisinstitution{Louisiana State University and Agricultural and Mechanical College}

% "Margins must be one-inch wide on all four sides of every page."
% "Margins must be the same on every page with no exceptions for wide tables
% and figures in landscape format."
%
% That second bullet is why no chapter-opening margin is set: LSU states the
% margin admits no exceptions at all, so a deeper opening would be out of spec
% rather than merely unasked-for.
\thesissetgeometry{left=1in, right=1in, top=1in, bottom=1in}

% A CHOICE, and an unusually wide one — LSU permits a single-spaced document:
% "Your document's narrative text may be either single- or double-spaced
% throughout."
%
% The profile takes double, which is the class default and the more common
% filing; \thesissetspacing{single} is equally in spec. "Throughout" is the
% operative word — LSU forbids mixing, and separately forbids anything between:
% "Documents should contain no spaces larger than a double space, except on the
% title page. Do not use half-spaces."
%
% LSU's per-element exceptions belong to the document rather than the class:
% the table of contents, lists and footnotes are always single-spaced even in a
% double-spaced document, and certain elements are always double-spaced even in
% a single-spaced one.
\thesissetspacing{double}

% --- Two rules this file cannot express, recorded so they are not lost --------
%
% PAGE NUMBERS: "Center all page numbers at the bottom of the page within the
% bottom one-inch margin, i.e., one-half inch from the bottom of the page" and
% "Format all page numbers in the same font and 12-point size as your text. Do
% not use boldface, italic, or ornamentation." The position matches the class's
% default foot folio; the half-inch placement does not follow from the profile
% fields, so a filer should check it.
%
% JUSTIFICATION: "Margins may be either left- or full-justified. Left
% justification is more flexible." A choice, so nothing is imposed.

% --- Accessibility -----------------------------------------------------------
% LSU publishes an accessibility rule and it is a CONDITION OF UPLOAD, so it is
% recorded — but it is a best-effort standard, not a technical one, and that
% distinction is the whole point of how it is encoded here.
%
% NOT set, deliberately: \thesisrequiretagging, \thesisrequiredocumentlanguage
% and \thesisrequirepdfstandard. LSU names no tagging, no document language, no
% WCAG level and no PDF standard. What it asks for is an effort — "to the best
% of your ability" — and turning that into a build-time conformance requirement
% would be this profile inventing a technical standard LSU never set. Compare
% University of Connecticut in the same round, which does name WCAG 2.1 and
% therefore does get the declarations.
\thesisaccessibilityrequirement{Accessible to visually impaired readers}
	{Make your document accessible to visually impaired readers to the best of
	 your ability. Refer to the LSU Digital Resources \& Content Accessibility
	 webpage and LaTeX's Guide to Produce Accessible PDF Files.}
\thesisaccessibilityrequirement{Condition of upload}
	{Do not upload your document until, to the best of your ability, you have
	 made the document accessible to the visually impaired. (Item 4 of the five
	 conditions the Graduate School sets before submission.)}

% --- Profile-local fields -----------------------------------------------------
% LSU's title page lists the author's prior degrees, one per line, under the
% name — "B.A., Purdue University, 2001 / M.A., University of Texas, 2004 /
% M.L.S., University of Virginia, 2010". The class has no slot for them.
% Defaults to empty, which is a real answer for a master's candidate with no
% prior graduate degree, so the lines simply do not print.
\newcommand{\thesis@lsu@priordegrees}{}
\newcommand{\lsupriordegrees}[1]{\renewcommand{\thesis@lsu@priordegrees}{#1}}

% --- Title page --------------------------------------------------------------
% "Meticulously follow the format shown on the sample title page, including
% placement of the three separate blocks of text, the use of double and single
% spaces, the words contained on each line, and the capitalization or
% lowercasing of every word."
%
% The three blocks, from the sample on handbook page 25:
%
%   (1) TITLE — "Format the title in solid capital letters, 16 points,
%       single-spaced, and centered on the first line below the top margin. The
%       rest of the text on the title page is 12 points." The sample sets it
%       bold; the prose does not mention weight, so the sample governs.
%
%   (2) A Dissertation
%       Submitted to the Graduate Faculty of the
%       Louisiana State University and
%       Agricultural and Mechanical College
%       in partial fulfillment of the
%       requirements for the degree of
%       Doctor of Philosophy
%       in
%       The Department of History
%
%   (3) by
%       Susan Mary Alford
%       B.A., Purdue University, 2001
%       M.A., University of Texas, 2004
%       M.L.S., University of Virginia, 2010
%       December 2020
%
% "Although the title page is page number i, it contains no page number."
% "On the final line, provide the month and year of graduation, e.g., May 2022."
% — so set \graddate to a bare "Month Year" for LSU; the class's default
% "May, 20XX" carries a comma LSU's example does not.
%
% \thesis@program supplies the department line, and the sample writes it with
% the article: "The Department of History". Set \gradprogram that way.
%
% "Use the version of your name that appears in official university records,
% which you will find on your Workday Unofficial Transcript. ... so that it
% will match on both your approval form and your title page."
\renewcommand{\thesistitlepage}{%
	\loadgeometry{nopagenumber}%
	\begin{titlepage}%
		\thispagestyle{empty}\singlespacing
		% "centered on the first line below the top margin": the geometry
		% already places the text block at 1in, so the title is the first thing
		% on the page with nothing above it. \topskip is zeroed for the reason
		% recorded in the Delaware profile — it would otherwise add 10pt before
		% the page's first box.
		\topskip=0pt\relax
		{\centering
			\thesis@tagtitle{Title}{%
				{\fontsize{16}{19}\selectfont\bfseries\MakeUppercase{\@title}\par}%
			}%
			% Block 2. The class's 12pt body size is the "rest of the text".
			\vfill
			A \thesis@DocNoun\par
			\vspace{\baselineskip}
			Submitted to the Graduate Faculty of the\par
			Louisiana State University and\par
			Agricultural and Mechanical College\par
			in partial fulfillment of the\par
			requirements for the degree of\par
			\thesis@degree\par
			\vspace{\baselineskip}
			in\par
			\vspace{\baselineskip}
			\thesis@program\par
			% Block 3.
			\vfill
			by\par
			\@author\par
			\ifx\thesis@lsu@priordegrees\@empty\else
				\thesis@lsu@priordegrees\par
			\fi
			\thesis@date\par
		\par}
	\end{titlepage}%
	\loadgeometry{normal}%
}

% --- Committee approval page --------------------------------------------------
% LSU prescribes none inside the document. Approval travels on a separate form
% that never becomes a page of the thesis: "a copy of the committee-signed
% Thesis/Dissertation Approval Form, which your department prepares and emails
% to gradsvcs@lsu.edu."
%
% That is stronger evidence than a closed list alone, and the closed list agrees
% with it. LSU's "Order of the Main Sections" runs Title Page, Copyright page,
% Dedication, Epigraph, Acknowledgments, Table of Contents, List of Tables,
% List of Figures, Nomenclature/Symbols/Acronyms, Abstract — with no approval,
% committee or signature page anywhere in the front matter.
%
% So the committee does sign at LSU; it signs a form that is emailed to the
% Graduate School, not a leaf of the filed document.
\thesisnoapprovalpage
